% ============================================================
\subsection{JPPipeline : énumération parallèle et statistiques}
\label{ann:jppipeline}
% ============================================================

% Brouillon rédigé par Claude le 2026-10-03, à reprendre dans la voix de Quentin. Commandes et
% comportements vérifiés ce jour-là en exécutant JPP (d = 2 et d = 3) dans un dossier temporaire.

JPP se compile et se lance depuis la racine du dépôt, avec \texttt{clingcon} dans le \texttt{PATH}, que fournit par exemple l'environnement conda du projet (\cref{ann:a-chaine}) :
{\small
\begin{verbatim}
javac -cp JPPipeline/lib/jpp.jar -d JPPipeline/bin \
      $(find JPPipeline/src/java_src/PSBN -name '*.java')
java -cp JPPipeline/bin:JPPipeline/lib/jpp.jar PSBN.Main \
      <dim> <lp_enumeratorFile> <setup_btreeFolder> <nSolutions> \
      <output_path> [includeWeightsAndTransitions]
\end{verbatim}
}
\begin{description}
\setlength{\itemsep}{0pt}
\item[\texttt{<dim>}] dimension $d$ des réseaux ;
\item[\texttt{<lp\_enumeratorFile>}] programme ASP qui énumère les réseaux (\cref{ann:asp}) ;
\item[\texttt{<setup\_btreeFolder>}] dossier contenant \texttt{setup\_btree\_<d>d.lp} et \texttt{Subdivision.lp}, en pratique \texttt{JPPipeline/src/asp} ;
\item[\texttt{<nSolutions>}] nombre maximal de solutions par vecteur de décomposition (0 pour toutes) ;
\item[\texttt{<output\_path>}] dossier où écrire le CSV ;
\item[{\texttt{[includeWeightsAndTransitions]}}] facultatif, 0 pour omettre les colonnes \texttt{f\_*} et \texttt{w\_*} (1 par défaut).
\end{description}
% src : Main.java (args[5] facultatif) ; PipeInput.java (noms des champs) ; EnumeratorPSBNFull.java et
% EnumeratorPSBN.java (setup_btree_<d>d.lp et Subdivision.lp lus dans setup_btreeFolder) ;
% conv a5891d3a 2026-06-12 (première version de cette doc)
Par exemple, la configuration \texttt{Main} de \texttt{.vscode/launch.json}, qu'on peut aussi lancer directement depuis VS Code, correspond à
{\small
\begin{verbatim}
java -cp JPPipeline/bin:JPPipeline/lib/jpp.jar PSBN.Main 2 \
      JPPipeline/src/asp/PSBN_undercontrol_taggedtree.lp JPPipeline/src/asp \
      0 JPPipeline/out 1
\end{verbatim}
}
% src : .vscode/launch.json ; .vscode/settings.json (sourcePaths, referencedLibraries, outputPath)

Le CSV, nommé \texttt{<dim>d\_<programme>\_output.csv}, s'écrit au fur et à mesure dans le dossier \texttt{<output\_path>}, et il vaut mieux le compresser ensuite avec \texttt{gzip}, puisque les pages HTML lisent directement les fichiers \texttt{.csv.gz}.
% src : Pipe.java (reset) ; DataStringLineCollectorProcessor.java ; sbn-csv.js
Pour suivre un long calcul, les options \texttt{-Djpp\_monitoring=true} et \texttt{-Djpp\_monitoring.stepTime=5}, placées avant \texttt{-cp} (elles figurent déjà dans \texttt{launch.json}), affichent toutes les cinq secondes l'avancement de chaque étape et la mémoire utilisée.
% src : .vscode/launch.json (vmArgs) ; conv a5891d3a 2026-06-05 et 2026-06-12 ; test 2026-10-03

L'argument \texttt{<nSolutions>} permet de ne demander à \texttt{clingcon} qu'un nombre limité de solutions par vecteur.
Comme \texttt{clingcon} est lancé avec une graine fixe, cette énumération partielle redonne les mêmes réseaux d'une exécution à l'autre.
% src : ClingconUtil.java (--seed=43, --rand-freq 1, --parallel 1) ; test 2026-10-03 (50 solutions par
% vecteur en d = 3 : mêmes 462 lignes sur deux exécutions, et avec le clingcon du stage)
Lorsque ce nombre est nul ou dépasse 1000, un garde-fou arrête l'énumération d'un vecteur à 1000 solutions et la relance en tranches, en fixant les premiers bits de la fonction du nœud 1 grâce à \texttt{Subdivision.lp}, pour limiter la mémoire de chaque appel ; le nombre de solutions demandé est alors réparti entre les tranches.
% src : EnumeratorPSBNFull.java (guardBinds, cap) ; Pipe.java (SUBDIVISION_THRESHOLD = 1000) ;
% EnumeratorPSBN.java (fixed_bits, quota userN / nombre de préfixes réalisables) ; Subdivision.lp (en-tête)
Ce seuil de 1000 solutions (\texttt{SUBDIVISION\_THRESHOLD}) et le nombre maximal d'énumérations non subdivisées menées en parallèle (10, par \texttt{setMaximumParallelTasks}) sont codés en dur dans \texttt{Pipe.java}, où il faut les modifier pour les adapter à la machine, les autres étapes utilisant déjà autant de fils d'exécution que la machine a de cœurs.
% src : Pipe.java (l. 34, 58, 106)
Limiter ainsi le nombre d'énumérations menées en parallèle évite d'engorger JPP.
% src : Quentin, 2026-10-03

Pour ajouter une statistique, on la calcule dans la méthode \texttt{process} de \texttt{Statifier.java}, qui reçoit chaque réseau et renvoie sa ligne de CSV, puis on ajoute son nom, à la bonne place, dans l'en-tête que construit la méthode \texttt{reset} de \texttt{Pipe.java} ; les statistiques qui portent sur les attracteurs sont, elles, calculées dans \texttt{SBN.toCsvRow}.
% src : Statifier.java ; Pipe.java (reset) ; SBN.java (toCsvRow)
